\begin{lemma}[Half-step identities for an optimizing step pair]
Let $F$ be a finite-dimensional real inner-product space, let
$1\le C_0<\infty$ and $0\le\kappa<\infty$, and let $v\in\Aclass_{m,C_0}$ in the
sense of \noderef{IsFitClass} be centered at $(0,1)$ with constant $C_0$ in the
sense of \noderef{IsCenteredAt} and satisfy the horizontal Jensen bound with
constant $\kappa$ at every time in the sense of \noderef{HasJensenBound}. Write
$f=W$ for the contact envelope of \noderef{ContactEnv}, $f^\ell=\Penv_\ell f$,
$g_\epsilon=\Menv_\epsilon v$ as in \noderef{JGap}. Let $\epsilon>0$ and let
$z=(x,t)$ with $t>\epsilon$, and let $(z^+,z^-)$ be an optimizing step pair of
size $\epsilon$ at $z$ in the sense of \noderef{IsStepPair}, with
$z^\pm=(x^\pm,t\pm\epsilon)$. Put
\begin{equation}\label{eq:dp-velocities}
  Q_+=\frac{x^+-x}{\epsilon},\qquad
  Q_-=\frac{x-x^-}{\epsilon}.
\end{equation}
Then
\begin{equation}\label{eq:dp-step-choice}
\begin{aligned}
  f^{2\epsilon}(z)&=f^\epsilon(z^+)-\frac{\epsilon}{2}|Q_+|^2,\\
  \Menv_{2\epsilon}v(z)&=g_\epsilon(z^-)+\frac{\epsilon}{2}|Q_-|^2,
\end{aligned}
\end{equation}
and
\[
  Q_+=\nabla_xf^{2\epsilon}(z),
\]
the gradient being that of \noderef{DPSpatialRegularity}. Although $z^-$ is the
backward step, $Q_-$ is oriented forward from $x^-$ to $x$; the spatial increment
for that step is $-\epsilon Q_-$.
\end{lemma}
\begin{proof}
Write $f=W=\Cenv_tv-\kappa^2t/2$ for the contact envelope of
\noderef{ContactEnv}, $z=(x,t)$, and recall $t>\epsilon>0$.

\emph{Step 0: hypotheses for the envelope lemmas.}
By \noderef{IsFitClass} the function $v$ is continuous on all of space-time and
is a viscosity solution of $F[v]=0$ there; by
\noderef{UniformSubquadratic}, applied with the centering point $(0,1)$, it has
subquadratic spatial growth on compact time intervals of $\R$ in the sense of
\noderef{SubquadraticGrowth}. By \noderef{BaryEnvSandwich} with $I=\R$, applied
at each parameter $\sigma>0$, and \noderef{ContactEnv},
\begin{equation}\label{eq:dpsi-sandwich}
  v(y,\tau)-\frac{\kappa^2\tau}{2}\le f(y,\tau)\le v(y,\tau)\qquad(\tau>0),
\end{equation}
so, exactly as for $v$, the function $f$ has subquadratic spatial growth on
compact subintervals of $(0,\infty)$: let $J\subseteq(0,\infty)$ be compact and
let $\varepsilon'>0$. If $J=\emptyset$ the required conclusion holds vacuously
with $M'=1$, since there is no $\tau\in J$; so assume $J\ne\emptyset$, and then
$\max J$ exists because $J$ is a nonempty compact set of reals. Apply
\noderef{SubquadraticGrowth} for $v$ with $J$ and $\varepsilon'/2$ to get
$M'>0$, then enlarge $M'$ so that
$\kappa^2\max J/2\le(\varepsilon'/2)M'^2$; for $|y|\ge M'$ and $\tau\in J$ the
two bounds of \eqref{eq:dpsi-sandwich} then give
$|f(y,\tau)|\le\varepsilon'|y|^2$. Also $f$ is continuous on
$\{(y,\tau):\tau>0\}$, because \noderef{BaryEnvDuality}, part (3), applied to
$v$ with $I=\R$, makes $(\sigma,\zeta)\mapsto\Cenv_\sigma v(\zeta)$ continuous at
every $(\sigma,\zeta)$ with $\sigma>0$, and \noderef{ContactEnv} composes this
with the continuous map sending $\zeta$ to its time coordinate.

Hence \noderef{QuadEnvAttained} applies to $f$ with $I=(0,\infty)$ and to $v$
with $I=\R$: in both cases the hypotheses of that lemma -- continuity on
$\Omega=\{(y,\tau):\tau\in I\}$ and subquadratic spatial growth on compact time
intervals of $I$ in the sense of \noderef{SubquadraticGrowth} -- are exactly the
two facts just established. Also \noderef{DPSpatialRegularity} applies to $v$
and $\kappa$, since its hypotheses on $v$, $C_0$, $\kappa$ and $f$ are the
hypotheses of the present lemma.

\emph{Step 1: the optimizers and the velocities.}
By \noderef{IsStepPair} there are $Y^+,Y^-\in F$ with
\begin{equation}\label{eq:dpsi-opt}
  \Penv_{2\epsilon}f(z)=f\bigl(Y^+,t+2\epsilon\bigr)
      -\frac{|Y^+-x|^2}{2\cdot2\epsilon},
  \qquad
  \Menv_{2\epsilon}v(z)=v\bigl(Y^-,t-2\epsilon\bigr)
      +\frac{|Y^--x|^2}{2\cdot2\epsilon},
\end{equation}
and $z^\pm=(x^\pm,t\pm\epsilon)$ with
$x^+=\tfrac12(x+Y^+)$ and $x^-=\tfrac12(x+Y^-)$. By
\eqref{eq:dp-velocities},
\[
  Q_+=\frac{x^+-x}{\epsilon}=\frac{Y^+-x}{2\epsilon},
  \qquad
  Q_-=\frac{x-x^-}{\epsilon}=\frac{x-Y^-}{2\epsilon},
\]
so
\begin{equation}\label{eq:dpsi-points}
  Y^+=x+2\epsilon Q_+,\quad x^+=x+\epsilon Q_+,
  \qquad
  Y^-=x-2\epsilon Q_-,\quad x^-=x-\epsilon Q_- .
\end{equation}
Consequently the quadratic costs split evenly along the straight paths:
\begin{equation}\label{eq:dpsi-costs}
  \frac{|Y^+-x|^2}{4\epsilon}=\epsilon|Q_+|^2,
  \quad
  \frac{|Y^+-x^+|^2}{2\epsilon}=\frac{|x^+-x|^2}{2\epsilon}
     =\frac{\epsilon|Q_+|^2}{2},
\end{equation}
and likewise with $Q_-$ in place of $Q_+$.

We also record once the elementary splitting inequality used in Steps 2 and 3:
for all $a,b\in F$,
\begin{equation}\label{eq:dpsi-split}
  \frac{|a+b|^2}{4\epsilon}\le\frac{|a|^2}{2\epsilon}+\frac{|b|^2}{2\epsilon},
\end{equation}
because expanding the inner products gives
$2|a|^2+2|b|^2-|a+b|^2=|a|^2-2\langle a,b\rangle+|b|^2=|a-b|^2\ge0$, and
dividing by $4\epsilon>0$ yields \eqref{eq:dpsi-split}. The slack is exactly
$|a-b|^2/(4\epsilon)$, so \eqref{eq:dpsi-split} is an equality precisely when
$a=b$.

\emph{Step 2: the forward identity.}
We first show
\begin{equation}\label{eq:dpsi-fwd-eq}
  \Penv_\epsilon f(z^+)=f\bigl(Y^+,t+2\epsilon\bigr)-\frac{\epsilon|Q_+|^2}{2}.
\end{equation}
Throughout this step we use \noderef{QuadEnvAttained}, part (1), for $f$ with
$I=(0,\infty)$, which Step 0 licenses. Two instances are needed, and both have
admissible time arguments because $t>\epsilon>0$:

(2a) at the center $z^+=(x^+,t+\epsilon)$ and the length $\epsilon$, where
$t+\epsilon\in(0,\infty)$ and $(t+\epsilon)+\epsilon=t+2\epsilon\in(0,\infty)$;

(2b) at the center $z=(x,t)$ and the length $2\epsilon$, where $t\in(0,\infty)$
and $t+2\epsilon\in(0,\infty)$.

Each instance supplies attainment: the supremum defining the envelope is equal
to the value at some competitor. That is all \noderef{QuadEnvAttained} gives,
and it is NOT the same as domination. The envelope is a supremum, and the
supremum of a set that is empty or unbounded above is a junk value, so an
equality with the value at one competitor does not by itself place the other
competitors below the envelope -- the inference runs the other way, since
attainment in the sense of a genuine maximum presupposes boundedness rather
than supplying it. The domination used below is therefore taken from part (1)
of \noderef{EnvelopeCompetitorBounds}, whose first clause states that
$\Penv_\ell f(\zeta)$ dominates every competitor value, applied at the same two
instances: at the centre $z^+$ with length $\epsilon$ for (2a), and at the
centre $z$ with length $2\epsilon$ for (2b). Its hypotheses are the continuity
and subquadratic growth of $f$ on $(0,\infty)$ established in Step 0, together
with the time memberships just checked. Only domination is needed; the
leastness half of the supremum is not used in this step.

For ``$\ge$'': by \eqref{eq:dpsi-points} the point $Y^+$ is a competitor in
\noderef{QuadEnvP} for $\Penv_\epsilon f(z^+)$, whose time argument is
$(t+\epsilon)+\epsilon=t+2\epsilon$, and its cost is
$|Y^+-x^+|^2/(2\epsilon)=\epsilon|Q_+|^2/2$ by \eqref{eq:dpsi-costs}. By
instance (2a) this competitor value is at most the envelope:
\begin{equation}\label{eq:dpsi-fwd-ge}
  \Penv_\epsilon f(z^+)\ge f\bigl(Y^+,t+2\epsilon\bigr)
    -\frac{\epsilon|Q_+|^2}{2}.
\end{equation}

For ``$\le$'': by instance (2a) the supremum defining $\Penv_\epsilon f(z^+)$ is
attained at some $y_1\in F$, that is,
\begin{equation}\label{eq:dpsi-fwd-max}
  \Penv_\epsilon f(z^+)=f\bigl(y_1,t+2\epsilon\bigr)
    -\frac{|y_1-x^+|^2}{2\epsilon}.
\end{equation}
The same point $y_1$ is a competitor in \noderef{QuadEnvP} for
$\Penv_{2\epsilon}f(z)$, whose time argument is also $t+2\epsilon$ and whose
cost is $|y_1-x|^2/(2\cdot2\epsilon)=|y_1-x|^2/(4\epsilon)$; by instance (2b)
that competitor value is at most the envelope, so
\begin{equation}\label{eq:dpsi-fwd-comp}
  f\bigl(y_1,t+2\epsilon\bigr)\le\Penv_{2\epsilon}f(z)
    +\frac{|y_1-x|^2}{4\epsilon}.
\end{equation}
Apply \eqref{eq:dpsi-split} with $a=y_1-x^+$ and $b=x^+-x$, whose sum is
$y_1-x$:
\begin{equation}\label{eq:dpsi-fwd-ineq}
  \frac{|y_1-x|^2}{4\epsilon}\le\frac{|y_1-x^+|^2}{2\epsilon}
    +\frac{|x^+-x|^2}{2\epsilon}.
\end{equation}
Substituting \eqref{eq:dpsi-fwd-comp} into \eqref{eq:dpsi-fwd-max} and then
using \eqref{eq:dpsi-fwd-ineq}, the two terms $|y_1-x^+|^2/(2\epsilon)$ cancel:
\[
  \Penv_\epsilon f(z^+)
   \le\Penv_{2\epsilon}f(z)+\frac{|y_1-x|^2}{4\epsilon}
     -\frac{|y_1-x^+|^2}{2\epsilon}
   \le\Penv_{2\epsilon}f(z)+\frac{|x^+-x|^2}{2\epsilon}
   =\Penv_{2\epsilon}f(z)+\frac{\epsilon|Q_+|^2}{2},
\]
the last equality by \eqref{eq:dpsi-costs}. By \eqref{eq:dpsi-opt} and
\eqref{eq:dpsi-costs}, $\Penv_{2\epsilon}f(z)
=f(Y^+,t+2\epsilon)-\epsilon|Q_+|^2$, so
\[
  \Penv_\epsilon f(z^+)\le f\bigl(Y^+,t+2\epsilon\bigr)-\epsilon|Q_+|^2
     +\frac{\epsilon|Q_+|^2}{2}
   =f\bigl(Y^+,t+2\epsilon\bigr)-\frac{\epsilon|Q_+|^2}{2}.
\]
With \eqref{eq:dpsi-fwd-ge} this proves \eqref{eq:dpsi-fwd-eq}. Note that no
property of the composed function $\Penv_\epsilon f$ is used anywhere above:
every supremum appearing is an envelope of $f$ itself.

Combining \eqref{eq:dpsi-opt}, \eqref{eq:dpsi-costs} and
\eqref{eq:dpsi-fwd-eq},
\[
  \Penv_{2\epsilon}f(z)=f\bigl(Y^+,t+2\epsilon\bigr)-\epsilon|Q_+|^2
   =\Bigl(\Penv_\epsilon f(z^+)+\frac{\epsilon|Q_+|^2}{2}\Bigr)
     -\epsilon|Q_+|^2
   =\Penv_\epsilon f(z^+)-\frac{\epsilon}{2}|Q_+|^2,
\]
which is the first line of \eqref{eq:dp-step-choice}.

\emph{Step 3: the backward identity.}
Symmetrically, we show
\begin{equation}\label{eq:dpsi-bwd-eq}
  \Menv_\epsilon v(z^-)=v\bigl(Y^-,t-2\epsilon\bigr)+\frac{\epsilon|Q_-|^2}{2}.
\end{equation}
Throughout this step we use \noderef{QuadEnvAttained}, part (1), for $v$ with
$I=\R$, which Step 0 licenses. Because $I=\R$, every time argument is
admissible with no sign condition -- which matters here, since $t>\epsilon$
leaves $t-2\epsilon$ possibly negative. Two instances are needed:

(3a) at the center $z^-=(x^-,t-\epsilon)$ and the length $\epsilon$, where
$t-\epsilon\in\R$ and $(t-\epsilon)-\epsilon=t-2\epsilon\in\R$;

(3b) at the center $z=(x,t)$ and the length $2\epsilon$, where $t\in\R$ and
$t-2\epsilon\in\R$.

Each instance supplies attainment: the infimum defining the envelope is equal to
the value at some competitor. As in Step 2 that is not domination, and for the
same reason with the inequalities reversed, since an infimum over a set that is
empty or unbounded below is equally the junk value. The domination used below is
part (2) of \noderef{EnvelopeCompetitorBounds}, whose first clause states that
$\Menv_\ell v(\zeta)$ is dominated by every competitor value, applied at the
centre $z^-$ with length $\epsilon$ for (3a) and at the centre $z$ with length
$2\epsilon$ for (3b); its hypotheses are the continuity and subquadratic growth
of $v$ on $\R$ established in Step 0. Only domination is needed here too.

For ``$\le$'': by \eqref{eq:dpsi-points} the point $Y^-$ is a competitor in
\noderef{QuadEnvM} for $\Menv_\epsilon v(z^-)$, whose time argument is
$(t-\epsilon)-\epsilon=t-2\epsilon$, with cost
$|Y^--x^-|^2/(2\epsilon)=\epsilon|Q_-|^2/2$ by \eqref{eq:dpsi-costs}. By
instance (3a) this competitor value is at least the envelope:
\begin{equation}\label{eq:dpsi-bwd-le}
  \Menv_\epsilon v(z^-)\le v\bigl(Y^-,t-2\epsilon\bigr)
    +\frac{\epsilon|Q_-|^2}{2}.
\end{equation}

For ``$\ge$'': by instance (3a) the infimum defining $\Menv_\epsilon v(z^-)$ is
attained at some $w_0\in F$, that is,
\begin{equation}\label{eq:dpsi-bwd-min}
  \Menv_\epsilon v(z^-)=v\bigl(w_0,t-2\epsilon\bigr)
    +\frac{|w_0-x^-|^2}{2\epsilon}.
\end{equation}
The same point $w_0$ is a competitor in \noderef{QuadEnvM} for
$\Menv_{2\epsilon}v(z)$, whose time argument is also $t-2\epsilon$ and whose
cost is $|w_0-x|^2/(2\cdot2\epsilon)=|w_0-x|^2/(4\epsilon)$; by instance (3b)
that competitor value is at least the envelope, so
\begin{equation}\label{eq:dpsi-bwd-comp}
  v\bigl(w_0,t-2\epsilon\bigr)\ge\Menv_{2\epsilon}v(z)
    -\frac{|w_0-x|^2}{4\epsilon}.
\end{equation}
Apply \eqref{eq:dpsi-split} with $a=w_0-x^-$ and $b=x^--x$, whose sum is
$w_0-x$:
\begin{equation}\label{eq:dpsi-bwd-ineq}
  \frac{|w_0-x|^2}{4\epsilon}\le\frac{|w_0-x^-|^2}{2\epsilon}
    +\frac{|x^--x|^2}{2\epsilon}.
\end{equation}
Substituting \eqref{eq:dpsi-bwd-comp} into \eqref{eq:dpsi-bwd-min} and then
using \eqref{eq:dpsi-bwd-ineq}, the two terms $|w_0-x^-|^2/(2\epsilon)$ cancel:
\[
  \Menv_\epsilon v(z^-)
   \ge\Menv_{2\epsilon}v(z)-\frac{|w_0-x|^2}{4\epsilon}
     +\frac{|w_0-x^-|^2}{2\epsilon}
   \ge\Menv_{2\epsilon}v(z)-\frac{|x^--x|^2}{2\epsilon}
   =\Menv_{2\epsilon}v(z)-\frac{\epsilon|Q_-|^2}{2},
\]
the last equality by \eqref{eq:dpsi-costs}. By \eqref{eq:dpsi-opt} and
\eqref{eq:dpsi-costs}, $\Menv_{2\epsilon}v(z)
=v(Y^-,t-2\epsilon)+\epsilon|Q_-|^2$, so
\[
  \Menv_\epsilon v(z^-)\ge v\bigl(Y^-,t-2\epsilon\bigr)+\epsilon|Q_-|^2
     -\frac{\epsilon|Q_-|^2}{2}
   =v\bigl(Y^-,t-2\epsilon\bigr)+\frac{\epsilon|Q_-|^2}{2}.
\]
With \eqref{eq:dpsi-bwd-le} this proves \eqref{eq:dpsi-bwd-eq}; again every
infimum used is an envelope of $v$ itself, never of a composed function, and
$w_0$ is not claimed to be unique. Combining with \eqref{eq:dpsi-opt} and
\eqref{eq:dpsi-costs},
\[
  \Menv_{2\epsilon}v(z)=v\bigl(Y^-,t-2\epsilon\bigr)+\epsilon|Q_-|^2
   =\Menv_\epsilon v(z^-)+\frac{\epsilon}{2}|Q_-|^2,
\]
which is the second line of \eqref{eq:dp-step-choice}.

\emph{Step 4: the forward velocity is the gradient.}
Apply \noderef{DPSpatialRegularity} to $v$ and $\kappa$ at the time $t>0$, the
length $2\epsilon>0$ and the point $x$. It supplies a point $y\in F$ which
attains $\Penv_{2\epsilon}f(x,t)$, is the \emph{unique} such point, and
satisfies
\[
  \text{$x'\mapsto\Penv_{2\epsilon}f(x',t)$ is differentiable at $x$ with
  gradient }\frac{y-x}{2\epsilon}.
\]
By \eqref{eq:dpsi-opt} the point $Y^+$ also attains $\Penv_{2\epsilon}f(x,t)$,
so uniqueness gives $Y^+=y$, and by \eqref{eq:dpsi-points}
$(y-x)/(2\epsilon)=Q_+$. Hence
$Q_+=\nabla_xf^{2\epsilon}(z)$, which is the last assertion.
\end{proof}
